% English translation of checked_pages/esq_119_055.tex; source: cote 119, site 55, PDF page 56, typescript page 53.
% AI-translated by GPT-6 Astra (OpenAI, ChatGPT Pro mode).
\scanpage{119/55}
\mspage{53}

the definition of the circle. It had impressed me by its
simplicity and obviousness, whereas the circle's property of ``perfect
roundness'' had previously seemed to me a mysterious reality
beyond words. It was at that moment, I believe, that I
glimpsed for the first time (without, of course, formulating it to myself
in these terms) the creative power of a ``good'' mathematical
definition, of a \uline{formulation} that describes the essence. Even today,
it seems that the fascination this power has exercised over
me has lost none of its force.

(\textsuperscript{3}) More generally, beyond so-called ``anabelian'' varieties
over fields of finite type, anabelian algebraic geometry
(as it emerged a few years ago) leads to a
description, in terms of profinite groups alone, of the category
of schemes of finite type over the absolute base $\mathbb{Q}$ (or
even $\mathbb{Z}$\frapparatus{V-FR-055-001}), and thereby, in principle, of the category of
arbitrary schemes (by suitable passages to the limit). It
is thus a construction that ``pretends'' to ignore rings
(such as $\mathbb{Q}$, algebras of finite type over $\mathbb{Q}$, etc.) and
the algebraic equations traditionally used to describe
schemes, by working directly with their étale topoi,
expressible in terms of systems of profinite groups. A
grain of salt, however: to be able to hope to reconstruct a
scheme (of finite type over $\mathbb{Q}$, say) from its étale topos,
which is a purely topological invariant, one should place oneself
not in the category of schemes (of finite type over $\mathbb{Q}$ in
this case), but in the one obtained from it by ``localization'',
making invertible the morphisms that are ``universal
homeomorphisms'', i.e. that are finite, radicial, and surjective. The
development of such a translation of a ``geometric world''
(namely that of schemes, scheme-theoretic multiplicities, etc.) into
terms of an ``algebraic world'' (that of profinite groups, and
systems of profinite groups describing suitable topoi (called ``étale''))
may be regarded as an ultimate outcome

\margin{55}
